\section{Introduction}
\label{sec:intro}

Multi-panel story image generation asks for a coherent visual sequence from a short narrative written as ordered scene blocks.
Unlike isolated text-to-image synthesis, each panel must be \emph{locally correct}---depicting the intended action and setting---and \emph{globally consistent}---keeping recurring characters recognizable across panels.
This dual requirement is difficult because independent per-scene generation often drifts in appearance, pronouns such as ``she'' or ``they'' must be mapped to the right subjects, and two-character scenes suffer from identity swapping or subject merging when a single visual anchor is forced onto both faces.

Our project builds a fully automated pipeline that separates textual story understanding from visual identity control.
An LLM-assisted planner produces strict structured metadata for characters and scenes; deterministic local code validates the plan, builds prompts, routes generation, and logs every decision.
The \textbf{submitted system} is a \emph{subject-count hybrid router}: stories with one distinct tagged entity use anchor-conditioned SDXL generation with IP-Adapter; stories with two or more entities use native StoryDiffusion, where front-loaded identity prompts initialize an internal identity bank.
This design follows a safety principle---single-anchor conditioning is skipped when two characters are equally important, rather than guessing one identity target.

After the first course presentation, we explored transformer-based diffusion (PixArt-$\alpha$) with story-level cross-scene attention and DreamBooth LoRA, as suggested for stronger long-range consistency.
That line is fully implemented and evaluated, but qualitative results remain below our hybrid baseline; we report it as an exploratory extension with explicit failure analysis.

\paragraph{Contributions.}
\begin{itemize}
  \item A reproducible structured prompt and identity framework spanning humans, animals, and robots, with subject-type-aware normalization.
  \item A subject-count hybrid generation policy combining IP-Adapter + SDXL for single-character stories and native StoryDiffusion for multi-character stories.
  \item Scene-consistency prompt assembly and CLIP-based panel selection for the single-character path.
  \item An exploratory PixArt-$\alpha$ DiT backend with LoRA and cross-scene blending, documented as a controlled negative comparison.
\end{itemize}

\begin{figure*}[!t]
  \centering
  \begin{tikzpicture}[
    node distance=0.55cm and 0.7cm,
    box/.style={draw, rounded corners, align=center, font=\small, minimum height=0.65cm, inner sep=3pt},
    arr/.style={-{Stealth[length=2mm]}, thick}
  ]
    \node[box] (in) {Scene-marked\\story text};
    \node[box, right=of in] (parse) {Parser};
    \node[box, right=of parse] (route) {Entity-count\\router};
    \node[box, below=0.9cm of route, xshift=-1.6cm] (single) {Single-character\\path};
    \node[box, below=0.9cm of route, xshift=1.6cm] (dual) {Multi-character\\path};
    \node[box, below=0.75cm of single] (llm) {LLM structured\\planner};
    \node[box, below=0.75cm of llm] (anchor) {Anchor bank\\+ IP-Adapter};
    \node[box, below=0.75cm of anchor] (sdxl) {SDXL-Turbo\\generation};
    \node[box, below=0.75cm of dual] (nat) {StoryDiffusion\\prompt adapter};
    \node[box, below=0.75cm of nat] (sd) {Native StoryDiffusion\\generation};
    \node[box, below=1.35cm of sdxl, xshift=1.6cm] (clip) {CLIP panel\\selection};
    \draw[arr] (in) -- (parse);
    \draw[arr] (parse) -- (route);
    \draw[arr] (route) -- node[left, font=\scriptsize]{1 entity} (single);
    \draw[arr] (route) -- node[right, font=\scriptsize]{$\geq$2} (dual);
    \draw[arr] (single) -- (llm);
    \draw[arr] (llm) -- (anchor);
    \draw[arr] (anchor) -- (sdxl);
    \draw[arr] (dual) -- (nat);
    \draw[arr] (nat) -- (sd);
    \draw[arr] (sdxl) -- (clip);
  \end{tikzpicture}
  \caption{Overview of the subject-count hybrid pipeline. Shared parsing and planning feed either anchor-conditioned SDXL (single entity) or native StoryDiffusion (multiple entities).}
  \label{fig:pipeline}
  % USER: replace with fig_pipeline_overview.pdf if a polished diagram is preferred
\end{figure*}
